\documentclass[10pt,twocolumn]{article}

\usepackage[utf8]{inputenc}
\usepackage[T1]{fontenc}
\usepackage[margin=1.9cm]{geometry}
\usepackage{amsmath,amssymb,amsthm}
\usepackage{booktabs}
\usepackage{graphicx}
\usepackage{xcolor}
\usepackage[expansion=false]{microtype}
\usepackage{titlesec}
\usepackage{caption}
\usepackage{hyperref}
\usepackage{enumitem}
\usepackage{abstract}

\captionsetup{font=small,labelfont=bf}
\hypersetup{colorlinks=true,linkcolor=blue!55!black,citecolor=blue!55!black,urlcolor=blue!55!black}

\titleformat{\section}{\normalfont\large\bfseries}{\thesection}{0.6em}{}
\titleformat{\subsection}{\normalfont\normalsize\bfseries}{\thesubsection}{0.6em}{}
\titlespacing*{\section}{0pt}{1.4ex plus 1ex minus .2ex}{0.9ex plus .2ex}

\newtheorem{proposition}{Proposition}
\newcommand{\gp}{\,{*}\,}            % geometric product
\newcommand{\real}{\mathbb{R}}
\newcommand{\Cl}{\mathrm{Cl}}

\title{\vspace{-1.2em}\bfseries Geometric Versor Memory and Geometric Attention:\\
A Clifford-Algebra Substrate for One-Shot Binding, Editing, and Order Encoding}

\author{%
Ignacio Mar\'ia Ozc\'ariz Arraiza\\[0.3em]
\small Doctor en Electr\'onica, Inform\'atica y Comunicaciones,\\[-0.3em]
\small Escuela de Ingenieros Industriales, Universidad de Extremadura\\[-0.2em]
\small Ingeniero Aeroespacial, Universidad Polit\'ecnica de Madrid\\[0.35em]
\small ORCID:~\href{https://orcid.org/0000-0002-9511-016X}{0000-0002-9511-016X}
\quad\href{mailto:anaciber@gmail.com}{anaciber@gmail.com}
}
\date{}

\begin{document}
\maketitle

\begin{abstract}
\noindent
Modern attention binds queries and keys by a symmetric dot product: it is blind
to order, and the knowledge it stores is expensive to edit. We study an
associative-memory and attention model that replaces this inner product with the
geometric (Clifford) product of unit \emph{versors}. Binding is a single
geometric product, recall is multiplication by the versor inverse, and---because
that inverse is exact---a stored fact is deleted exactly, with no retraining. We
measure the model on small algebras $\Cl(n,0)$, $n\le 8$, under a reproducible
synthetic protocol, and report both what it delivers and what it does not. A
single field has low, sub-linearly scaling capacity dominated by superposition
cross-talk (E1), which a sparse addressable arrangement restores to linear (E2);
the pure model generalizes by robustness and exact analogy but needs a softmax
readout to interpolate (E3--E4); and against classical one-shot and trained
baselines it loses in raw capacity and noise robustness, its measured edge being
storage density and composable algebraic structure (E0). Two trained-model
results follow. Versor position binding is \emph{provably} orthogonal, generalizes
rotary position embeddings (RoPE), and matches them once the frequency spectrum
is designed (E6). And, as our flagship result, under \emph{sequential} knowledge
editing a frozen-model versor memory surpasses a faithful rank-one (ROME)
baseline on all three canonical metrics at 60 accumulated edits---efficacy
$1.00$ vs.\ $0.82$, paraphrase $0.70$ vs.\ $0.43$, locality $0.98$ vs.\
$0.92$---with $O(1)$ edit cost and exact deletion (E7). We conclude that the
model is best understood not as a ``one-shot large language model'' but as an
edit-friendly, order-aware \emph{memory--binding substrate} for attention, in
which the attention inner product is the geometric product. All results are
small-model and synthetic; scaling to pretrained language models and standard
benchmarks is left to future work.
\end{abstract}

\section{Introduction}
Modern sequence models are dominated by the transformer, whose core operation is
scaled dot-product attention, $\mathrm{softmax}(QK^\top/\sqrt{d})\,V$~\cite{vaswani2017}. Two
recurring weaknesses of this paradigm are (a) the cost and instability of
\emph{editing} stored knowledge---adding or removing a fact typically requires
re-training or fine-tuning---and (b) the reliance on bolted-on positional
encodings because the dot product is symmetric and therefore blind to order.
Vector Symbolic Architectures (VSA) and Holographic Reduced Representations
(HRR)~\cite{plate1995,kanerva2009,smolensky1990,kleyko2022} offer an alternative:
represent items as high-dimensional vectors and bind them
by a multiplication-like operation, enabling one-shot storage and algebraic
composition. Geometric (Clifford) algebra~\cite{hestenes1984,dorst2007} refines
this idea by giving the binding
operation a geometric meaning---rotations and reflections---and by separating an
inner (symmetric) and outer (antisymmetric) part.

This paper formalizes and \emph{measures} a geometric-algebra memory in which
the state is a multivector, items are unit \emph{versors}, binding is the
geometric product, and recall is multiplication by the versor inverse. We then
ask the empirical questions that decide whether such a model is viable:
how much can a field store; does sparsity restore scaling; does it generalize;
and what does a geometric attention layer add over a standard one. We answer each
with reproducible experiments on $\Cl(n,0)$ for $n\le 8$. Our contributions are
deliberately scoped: they hold in a small-algebra, synthetic-data regime and are
intended as a proof-of-concept and a precise statement of which properties are
genuine and which are limited.

\paragraph{Contributions.} Concretely, we contribute:
(1)~a versor associative memory with exact, invertible binding, and an analysis
(E1) showing that the \emph{conditioning} of the unbinding map---not grade
richness---governs single-field capacity;
(2)~a sparse addressable construction (E2) that removes the single-field ceiling
and restores linear capacity scaling;
(3)~a geometric-attention layer (E3--E4) that recovers interpolation through
softmax while \emph{uniquely} retaining an order/orientation dimension the dot
product discards;
(4)~an honest baseline study (E0) locating the model's advantage in storage
density and composable algebraic structure, not raw capacity or robustness;
(5)~a proof, with a trained-transformer demonstration (E6), that versor position
binding generalizes RoPE, the decisive variable being the frequency spectrum;
and
(6)~a flagship sequential-editing result (E7) in which a frozen-model versor
memory surpasses a faithful rank-one ROME baseline on all three canonical
metrics at 60 edits.

\section{Background and Model}

\subsection{Geometric algebra $\Cl(n,0)$}
We work in the Clifford algebra generated by $n$ orthonormal vectors
$e_1,\dots,e_n$ with $e_ie_j+e_je_i = 2\delta_{ij}$ and positive signature. The
algebra has $2^n$ basis blades graded by the number of generators they contain
(scalar, vectors, bivectors, \dots, pseudoscalar). A multivector $a$ is a real
linear combination of blades. The \emph{geometric product} decomposes into a
symmetric inner part and an antisymmetric outer part,
\begin{equation}
a \gp b \;=\; a\cdot b \;+\; a\wedge b ,\qquad
a\cdot b = b\cdot a,\quad a\wedge b = -\,b\wedge a .
\label{eq:gp}
\end{equation}
For the four-vector case $\Cl(4,0)$ used in the original neuro-quantum
formulation, a single geometric product expands to $16\times16=256$ scalar
multiplications.

\subsection{Versors and exact inversion}
A \emph{versor} is a product of invertible vectors, $V=v_1\gp v_2\gp\cdots\gp v_k$.
Versors form a group; the reverse $\widetilde V$ (blade order reversed, inducing
the sign $(-1)^{g(g-1)/2}$ per grade $g$) yields the inverse
\begin{equation}
V^{-1} \;=\; \frac{\widetilde V}{V\gp\widetilde V},\qquad V\gp\widetilde V \in \real .
\label{eq:inv}
\end{equation}
We verified numerically that genuine versors invert to machine precision
($\max$ error $\sim 10^{-16}$ for $n\le 8$), whereas an \emph{arbitrary}
multivector does not, because $X\gp\widetilde X$ is then not scalar. This is the
first design constraint: keys and operators must be versors for recall to be
exact.

\subsection{Versor associative memory}
Given key/value versor pairs $(k_i,v_i)$, we store by binding and superposing,
and recall by unbinding with the key inverse:
\begin{align}
M &= \sum_{i=1}^{K} k_i \gp v_i, \label{eq:store}\\
\widehat v_q &= k_q^{-1}\gp M
   = v_q + \underbrace{\sum_{i\neq q} k_q^{-1}\gp k_i \gp v_i}_{\text{cross-talk}} .
\label{eq:recall}
\end{align}
A non-linear \emph{clean-up}---nearest stored value by cosine---resolves the
noisy readout to a discrete item. \emph{Editing} is exact and local: a stored
pair is removed by subtracting $k_i\gp v_i$, and a versor operator is undone by
multiplying with its inverse \eqref{eq:inv}.

\section{Experiments}
All experiments use $\Cl(n,0)$, synthetic random versors, and the protocol of
Eqs.~\eqref{eq:store}--\eqref{eq:recall}. Code and per-run JSON are released with
the paper. Reported numbers are means over repeated trials (4--200 depending on
cost); ``capacity'' is the largest number of associations recalled at a stated
clean-up accuracy threshold.

\subsection{E1: Single-field capacity and the role of encoding}
We measure capacity at $\ge 95\%$ clean-up accuracy for four code families:
grade-2 versors (\textsc{v2}, the original choice), mixed-grade versors
(\textsc{v\_n}, product of $n$ vectors), even rotors (\textsc{rot}), and
random full-spectrum unit multivectors with an exact matrix inverse
(\textsc{full}). Table~\ref{tab:enc} and Fig.~\ref{fig:cap} show that mixed-grade
versors roughly double the capacity of grade-2 versors, while \textsc{full}
codes \emph{under}-perform despite carrying a complete grade spectrum.

\begin{table}[t]
\centering
\caption{Single-field capacity at $\ge 95\%$ accuracy, by code family.
Mixed-grade versors (\textsc{v\_n}) dominate; full-spectrum codes are worst.}
\label{tab:enc}
\small
\begin{tabular}{lrrrrr}
\toprule
$n$ & dim & \textsc{v2} & \textsc{v\_n} & \textsc{rot} & \textsc{full}\\
\midrule
4 & 16  & 1 & 2 & 3 & 2\\
5 & 32  & 1 & 2 & 3 & 3\\
6 & 64  & 3 & 5 & 5 & 1\\
7 & 128 & 5 & 7 & 5 & 2\\
\bottomrule
\end{tabular}
\end{table}

\begin{figure}[t]
\centering
\includegraphics[width=\columnwidth]{fig_capacity.pdf}
\caption{Single-field capacity per unit dimension falls with algebra dimension
for every code family (E1): the single-field capacity is sub-linear. Mixed-grade
versors and rotors lead; full-spectrum codes, worst-conditioned at unbinding,
collapse at high dimension.}
\label{fig:cap}
\end{figure}

The mechanism is the \emph{conditioning of the unbinding map}. Writing the
geometric product $k^{-1}\gp(\cdot)$ as a linear operator on $\real^{2^n}$, its
condition number determines how recall amplifies cross-talk noise. We measured
(over 50 keys, $n=7$) a median condition number of exactly $1.000$ for versor
unbinding versus $\approx 19.5$ (90th pct.\ $35.6$) for full-multivector
unbinding. A versor unbinding is norm-preserving and does not amplify noise;
a full multivector's does. Thus the versor constraint, needed for exact
inversion, is simultaneously what keeps recall well-conditioned. A structural
observation explains why even mixed-grade versors fall short of an ideal code:
a versor occupies blades of a single parity (all even or all odd), so it is
never full-spectrum.

\paragraph{Finding.} Part of the capacity limit is the encoding (curable: prefer
mixed-grade versors), and part is intrinsic: even the best versor code scales
sub-linearly, with capacity/dimension falling from $\sim0.19$ (dim 16) to
$\sim0.05$ (dim 256).

\subsection{E2: Sparse addressable memory restores linear scaling}
Instead of one saturated field, we use $M$ small fields (``slots''), each holding
a few associations below its cross-talk ceiling, plus an addressing function that
routes a key to a slot. We compare a deterministic key-hash (O(1), content-blind)
and a content-addressable nearest-anchor scheme (O(M)). Using $\Cl(6,0)$
(dim 64) mixed-grade versors as the building block, total capacity at
$\ge 90\%$ accuracy scales \emph{linearly} with $M$
(Table~\ref{tab:sparse}, Fig.~\ref{fig:sparse}); we verified linearity to
$M=256$ (768 associations at $94\%$). The single-field ceiling ($\sim 4$) is
removed; capacity becomes a design parameter.

\begin{table}[t]
\centering
\caption{Sparse memory: total capacity at $\ge 90\%$ vs.\ number of slots $M$
(building block $\Cl(6,0)$, dim 64).}
\label{tab:sparse}
\small
\begin{tabular}{rrr}
\toprule
$M$ slots & hash (O(1)) & anchor (O(M))\\
\midrule
1  & 4   & 6\\
2  & 8   & 8\\
4  & 16  & 16\\
8  & 32  & 24\\
16 & 48  & 48\\
32 & 128 & 64\\
64 & 192 & 128\\
\bottomrule
\end{tabular}
\end{table}

\begin{figure}[t]
\centering
\includegraphics[width=\columnwidth]{fig_sparse.pdf}
\caption{Total capacity scales linearly in the number of slots $M$ (E2),
removing the single-field ceiling. Hash addressing is $O(1)$ but load-imbalanced;
content-addressable anchor addressing is smoother at $O(M)$ cost per access.}
\label{fig:sparse}
\end{figure}

\paragraph{Honest costs.} Hash addressing follows balls-in-bins: at mean load 3,
the most-loaded slot held up to 10 items, the source of the residual $\sim5\%$
error; consistent hashing or rebalancing mitigates this at the cost of state.
Anchor addressing is content-addressable but costs O(M) per access
($\approx 12\,\mu$s at $M{=}128$ vs.\ $0.6\,\mu$s for hash).

\subsection{E3: Generalization---noise, analogy, interpolation}
We probe three increasingly demanding forms of generalization on $\Cl(6,0)$.

\emph{T1 (noise robustness).} Querying with a corrupted key
$k_i+\eta$ (renormalized), recall accuracy degrades gracefully: $0.98$ at
$\eta{=}0$, $0.90$ at $\eta{=}0.2$, $0.58$ at $\eta{=}0.5$. A hash table is exact
at $\eta{=}0$ and $0$ otherwise; the associative memory tolerates corruption.
This claim must, however, be read against the right baseline: as E0 shows
(Sec.~\ref{sec:e0}), at \emph{equal load} a one-shot Hebbian associator and an
explicit-storage Hopfield memory are \emph{more} noise-robust than the versor
field (at $K{=}6$, $\eta{=}0.2$: $1.00$ and $1.00$ vs.\ $0.67$). Graceful
degradation is an advantage over exact lookup, not over classical associative
memories.

\emph{T3 (structured analogy).} Encoding concepts as superpositions of role
versors ($\textit{king}=\textit{royal}+\textit{male}$, etc.), the vector analogy
$\textit{king}-\textit{man}+\textit{woman}$ retrieves \textit{queen} with
accuracy $1.00$ across $n\in\{4,6,8\}$. This is exact symbolic generalization.

\emph{T2 (interpolation).} Storing a smooth scalar field $f(t)$ keyed by versors
along a path and querying at unseen midpoints, the linear-readout memory
\emph{fails}: mean absolute error $0.795$, exceeding the signal scale
$0.707$ (linear interpolation reference: $0.100$). Diagnosis shows the
key-similarity kernel along the path \emph{is} smooth and localized; the failure
is the linear readout, which averages over an over-broad tail. Replacing the
raw-cosine average with a softmax-weighted average over the same similarities
reduces error monotonically with temperature $T$ (from $0.53$ at the linear
baseline to $0.078$ at $T{=}0.02$), i.e.\ below the signal scale: genuine
interpolation.

\paragraph{Finding.} The model generalizes by robustness and exact analogy but
not by smooth interpolation, and the missing ingredient is precisely the softmax
non-linearity, not the geometric similarity metric.

\subsection{E4: Geometric attention}
Motivated by E3, we build a single attention layer in which the score is derived
from the geometric product rather than the dot product, then apply softmax and a
weighted value sum. We test three properties.

\emph{(A) Interpolation recovered.} With softmax readout, both dot-product and
geometric-score attention interpolate the E3-T2 task (error $\approx 0.11 \ll
0.71$). The non-linearity, not the score type, restores interpolation.

\emph{(B) Analogy preserved.} Retrieving the $\textit{king}-\textit{man}+
\textit{woman}$ analogy through attention preserves accuracy
($0.96$ at $n{=}4$, $1.00$ at $n\in\{6,8\}$).

\emph{(C) Order/orientation encoded.} Because
$a\cdot b=b\cdot a$, dot-product scoring cannot distinguish the ordered pair
$(a,b)$ from $(b,a)$: over 2000 pairs the discrimination margin
$|\,\mathrm{score}(a,b)-\mathrm{score}(b,a)\,|$ is exactly $0.00$. The geometric
product retains the antisymmetric wedge $a\wedge b=-\,b\wedge a$, giving a margin
of $1.69$. This is an information-theoretic distinction, not a statistical one.

\paragraph{Caveat.} On a realistic next-token sequence-retrieval task with
distinct value vectors, \emph{standard} dot-product attention outperformed the
single-field geometric binding ($1.00$ vs.\ $0.54$): the practical bottleneck is
field capacity (E1), which classical key--value attention sidesteps. The
defensible claim is therefore narrow and precise: the geometric product
\emph{retains} an order/orientation dimension that the dot product discards, which
helps exactly when a task needs relational or oriented structure (graphs,
geometry, syntax, transformations), not in pure key--value lookup.

\subsection{E0: Trained and explicit-storage baselines}
\label{sec:e0}
As a control across all preceding tasks, we run the identical protocol
(store $K$ key--value pairs at $d{=}64$; nearest-value clean-up) on three
classical systems with \emph{stated budgets}: a one-shot Hebbian linear
associator $W=\sum_i v_i k_i^\top$ ($d^2$ floats); a modern Hopfield memory with
explicit pattern storage and softmax retrieval~\cite{ramsauer2020,hopfield1982}
($2Kd$ floats, the memory mechanism of attention); and a two-layer MLP trained by
Adam under a fixed 3{,}000-step budget ($8{,}320$ parameters; the
``weights-as-memory'' regime of language models). Table~\ref{tab:e0} reports
capacity, storage, density, and write/delete properties; noise robustness is
measured at the \emph{same} load $K{=}6$ for every system.

\begin{table}[t]
\centering
\caption{E0: identical task, stated budgets ($d{=}64$, threshold $90\%$).
``$\ge$'' marks the scan ceiling; Hopfield was verified interference-free to
$K{=}2048$. Noise column: accuracy at fixed load $K{=}6$, $\eta{=}0.2$.}
\label{tab:e0}
\small
\setlength{\tabcolsep}{3.4pt}
\begin{tabular}{lrrrllr}
\toprule
system & cap@90 & storage & dens. & write & delete & noise\\
\midrule
versor      & 6          & 64      & \textbf{0.094} & 1-shot & exact   & 0.67\\
Hebbian     & 180        & 4{,}096 & 0.044          & 1-shot & exact   & \textbf{1.00}\\
Hopfield    & $\ge$2048  & $2Kd$   & 0.008          & append & exact   & \textbf{1.00}\\
MLP         & $\ge$256   & 8{,}320 & 0.031          & 3k steps & retrain & 0.94\\
\bottomrule
\end{tabular}
\end{table}

Three honest conclusions follow. First, in \emph{raw capacity} the classical
baselines dominate: explicit storage has no interference ceiling at all, and even
the Hebbian associator stores $30\times$ more than the versor field at this
dimension. Second, in \emph{equal-load noise robustness} the versor field is
mid-pack, not best (cf.\ the T1 qualification above); the gradient-trained MLP,
by contrast, is robust at low load but collapses with load ($0.11$ at $K{=}128$,
$\eta{=}0.2$). Third, one-shot writing and exact deletion are \emph{not} unique
to the versor memory---the Hebbian and explicit-Hopfield systems share both; the
MLP has neither (writing costs epochs, deletion costs retraining). What survives
as distinctive, by measurement, is (a) \emph{density}: the versor field is the
most compact fixed-size superposed store, $2$--$12\times$ the alternatives in
associations per float, relevant wherever state must be small and fixed (embedded
or FPGA settings); and (b) \emph{algebraic structure}: it is the only system of
the four whose binding is multiplicative, invertible, and composable---the basis
of the exact analogy (E3), operator composition, and order encoding (E4) that
neither the outer product nor dot-product attention provides. E0 thus sharpens
the value proposition: not raw capacity or robustness, but a compact fixed-size
state with composable algebraic structure.

\subsection{E6: Versor position binding in trained transformers}
\label{sec:e6}
We test the order-encoding claim inside a \emph{trained} model: a 2-layer,
1-head transformer ($d{=}32{=}\dim\Cl(5,0)$, JAX/Adam, 2 seeds per condition),
where only the position mechanism varies: no positions; additive sinusoidal PE;
rotary position embeddings (RoPE)~\cite{su2021roformer}; binding by powers of a
\emph{random} unit versor applied to $q,k$; and a \emph{spectrum-matched} versor
(designed frequencies in random, non-axis-aligned planes). Two theoretical facts,
verified at $10^{-16}$, motivate the comparison: left-multiplication by a unit
versor is \emph{orthogonal}, hence
$\langle P_t q, P_s k\rangle=\langle q, P_{s-t}k\rangle$---the relative-position
property of RoPE---and RoPE is exactly the special case of commuting,
axis-aligned 2-D rotors. Versor binding therefore \emph{generalizes} RoPE.

Empirically: on position retrieval (output the token at a queried position,
$L{=}10$), the no-position control sits at chance ($0.10\approx 1/L$,
confirming the task is impossible without positions) while sinusoidal, RoPE,
and versor binding all reach $1.00$. On pairwise order (does $a$ precede $b$?)
with length extrapolation (train $L\in[6,12]$, test $L\in\{16,20,24\}$), the
\emph{random} versor fails ($0.70$ in-distribution, near-chance extrapolated).
The diagnosis is spectral, not algebraic: the random versor's slowest rotation
channel has period ${\sim}6$ positions---it wraps \emph{inside} the training
context---whereas RoPE's slowest channel has period ${\sim}3.5{\times}10^{4}$.
A versor with RoPE's frequency spectrum in random non-commuting planes recovers
full parity: in-distribution $0.947$ vs.\ $0.954$, extrapolation
$0.78/0.69/0.60$ vs.\ $0.81/0.69/0.58$ (statistically indistinguishable).

\paragraph{Finding.} Versor binding is a working position mechanism in trained
transformers and is the mathematical generalization of RoPE; the decisive design
variable is the \emph{frequency spectrum}, not the algebraic form. The result is
parity with RoPE, not superiority; and a random versor is a provably poor
positional encoder.

\subsection{E7: Sequential knowledge editing (flagship result)}
\label{sec:e7}
Inverse-based editing is the capability where the substrate has no classical
counterpart among weight-editing methods, so we evaluate it head-to-head in a
miniature of the standard protocol~\cite{meng2022rome}. A base LM (2 layers,
$d{=}32$) is trained from scratch to memorize 120 facts
$(s,r)\!\to\!o$, each expressed through 4 paraphrase templates (3 in training,
1 held out; base accuracy $1.00$ train / $0.94$ held-out). We then apply up to
60 \emph{sequential} edits and track the three canonical metrics: efficacy
(edited facts answer the new object), paraphrase (same, on the held-out
template), and locality (unedited facts survive). Editors: naive fine-tuning;
a faithful rank-one ROME implementation (key $k^*$ from the FFN hidden state,
value $v^*$ optimized for the new object, covariance-whitened update
$W_2 \mathrel{+}= (v^*-W_2k^*)(C^{-1}k^*)^\top/(k^{*\top}C^{-1}k^*)$)
\cite{meng2022rome}; and the \emph{versor module}: base model frozen, each edit
written one-shot as a versor binding into the sparse fields of E2
(load $\le$ 2--3 per slot), retrieval gated by cosine similarity of the frozen
model's representation (threshold calibrated on training data: within-fact
$0.95$ vs.\ between-fact $0.87$).

\begin{figure*}[t]
\centering
\includegraphics[width=0.92\textwidth]{fig_editing.pdf}
\caption{Sequential knowledge editing (E7), our flagship result. As edits
accumulate, the rank-one ROME baseline degrades on efficacy and locality because
its deltas share one weight matrix and interfere; the frozen-model versor memory
stays flat because each edit lives in its own slot and the model is never
touched. At 60 edits the versor module (load 2, $\star$) dominates ROME on all
three canonical metrics, with $O(1)$ edit cost and exact deletion.}
\label{fig:e7}
\end{figure*}

\begin{table}[t]
\centering
\caption{E7: sequential editing, metrics vs.\ number of accumulated edits $k$
(efficacy / paraphrase / locality). Naive FT collapses (locality $0.04$ at
$k{=}20$, omitted). Bottom row: versor with slot load 2.}
\label{tab:e7}
\small
\setlength{\tabcolsep}{4.2pt}
\begin{tabular}{lccc}
\toprule
$k$ & ROME & versor (load 3) \\
\midrule
10 & 1.00 / 0.50 / 1.00 & 0.90 / 0.80 / 1.00\\
20 & 0.95 / 0.40 / 1.00 & 0.80 / 0.60 / 1.00\\
40 & 0.88 / 0.35 / 0.98 & 0.85 / 0.58 / 0.98\\
60 & 0.82 / 0.43 / 0.92 & 0.88 / 0.62 / 0.98\\
\midrule
60 & --- & \textbf{1.00 / 0.70 / 0.98} (load 2)\\
\bottomrule
\end{tabular}
\end{table}

Table~\ref{tab:e7} and Fig.~\ref{fig:e7} show the crossover. At few edits
($k\le 20$) ROME matches or beats the versor module on efficacy. As edits
accumulate, ROME degrades monotonically (efficacy $1.00\!\to\!0.82$, locality
$1.00\!\to\!0.92$)---its rank-one deltas share the same weight matrix and
interfere, the documented failure mode of sequential weight editing---while the
versor module stays flat, because each edit lives in its own slot and the model
is never touched. Error anatomy at $k{=}60$ closes the loop with E2: the versor
gate makes \emph{zero} threshold misses and \emph{zero} wrong-entry matches; the
residual efficacy loss ($0.88$) is exactly the slot cross-talk that E2's capacity
curve predicts at load 3. Reducing the load to 2 yields \textbf{efficacy $1.00$,
paraphrase $0.70$, locality $0.98$ at 60 edits}, against ROME's $0.82/0.43/0.92$.
Edit cost is $O(1)$: all 60 versor edits complete in under one second, versus a
per-edit $v^*$ optimization for ROME and gradient descent for FT. Deleting all
60 edits restores the base model exactly ($1.000$); honestly, ROME's stored
deltas are also subtractable ($1.000$ here), but the versor module never
modifies the model in the first place, so restoration is trivial by
construction, while FT has no mechanism ($0.033$).

\paragraph{Finding.} The structural asymmetry---weight edits interfere,
slot-isolated algebraic edits do not---produces the first measured
\emph{advantage over a state-of-the-art reference method}: at 60 sequential
edits the versor module dominates ROME on all three canonical metrics. The
caveats are stated in Sec.~\ref{sec:limits}; the asymmetry itself is what we
expect to persist at scale, and validating it on a pretrained LLM with
CounterFact/zsRE is the concrete next step.

\section{Discussion}
The seven experiments support a single reframing, with E0 fixing its honest
boundary and E6--E7 supplying trained-model evidence. The geometric-versor model is
\emph{not} a ``one-shot large language model'': its single-field capacity is
small and it does not, by itself, interpolate. It is, however, a faithful
\emph{memory--binding substrate} for attention. Concretely, consider a
transformer-style layer
$\mathrm{softmax}\!\big(s(Q,K)/T\big)\,V$ in which the score $s$ is read from the
geometric product $Q\gp K$ rather than from $Q K^\top$. Relative to the standard
layer this substrate adds three capabilities demonstrated above:
\begin{enumerate}[leftmargin=1.3em,itemsep=2pt,topsep=2pt]
\item \textbf{Exact structural binding} via the geometric product, with
unbinding by the versor inverse---preserving the analogy behavior of E3-T3/E4-B.
\item \textbf{Inverse-based editing}: a stored fact or an applied operator is
removed exactly by multiplying with its inverse \eqref{eq:inv}, without
re-training---realized in E7 as a frozen-model editor that surpasses a rank-one
ROME baseline under sequential editing.
\item \textbf{Intrinsic order/orientation encoding} from the antisymmetric part
of the product (E4-C), which the symmetric dot product cannot represent---realized
in E6 as a position-binding mechanism that provably generalizes RoPE.
\end{enumerate}
The non-linearity that the pure algebraic model lacked---softmax with a
controllable temperature---is the same one the transformer already uses, and it
is cheap to add (E3/E4-A). The well-conditioned versor unbinding (E1) and the
sparse addressable arrangement (E2) make the substrate numerically stable and
scalable in capacity.

\section{Limitations}
\label{sec:limits}
We state these plainly, as they bound every claim above.
\textbf{Scale.} Memory results are for $\Cl(n,0)$ with $n\le 8$ (dimension
$\le256$); the $2^n$ blade count makes the dense geometric product expensive, and
we have not measured behavior at the dimensions where VSA/HRR are normally
deployed ($10^3$--$10^4$). The trained models of E6--E7 are deliberately small
(2 layers, $d{=}32$).
\textbf{Data.} Keys, values, the interpolation target, and the E7 knowledge base
(120 synthetic facts, 4 templates) are synthetic; no natural-language or
perceptual data is used.
\textbf{Baselines.} E0 adds one-shot Hebbian, explicit-storage modern-Hopfield,
and small gradient-trained MLP baselines---and the versor memory loses to them
in raw capacity and equal-load robustness. E6--E7 add trained transformers,
RoPE, and a faithful rank-one ROME implementation, but \emph{not} pretrained
LLMs or the standard CounterFact/zsRE benchmarks; the E7 advantage is
demonstrated in a miniature of that protocol.
\textbf{Editing scope.} The E7 gate is calibrated on training-data
representations (within-fact $0.95$ vs.\ between-fact $0.87$); paraphrase
generalization is bounded by the frozen base model's representation invariance
(ceiling ${\sim}0.70$). At few edits ($k\le20$) ROME matches or exceeds the
versor module; the advantage emerges under accumulation. Exact deletion is also
available to ROME if deltas are stored.
\textbf{Spectrum design.} E6 shows a \emph{random} versor is a poor positional
encoder (its spectrum wraps within ${\sim}6$ positions); position-binding
claims hold only for designed spectra, where the result is parity with RoPE,
not superiority.
\textbf{Capacity.} Even with the best encoding, single-field capacity scales
sub-linearly; the linear scaling of E2 is in the number of fields, with the
addressing trade-offs noted.
\textbf{Order claim.} E4-C establishes that order information is \emph{available}
in the geometric score, not that it improves end-task accuracy in general; E4's
caveat shows a case where it does not.

\section{Related work and conclusion}
The model sits at the intersection of VSA/HRR~\cite{plate1995,kanerva2009,kleyko2022},
modern associative memories, and
geometric-algebra neural networks~\cite{brandstetter2022,ruhe2023,brehmer2023}. A key
nearby result is the modern Hopfield network~\cite{ramsauer2020,hopfield1982}, whose
continuous update rule was shown to be equivalent to transformer
attention~\cite{vaswani2017} and which stores exponentially many patterns with
one-step retrieval; our work differs in using the geometric product---not the dot
product---as the similarity kernel, and in foregrounding inverse-based editing and
order encoding. The Geometric Algebra Transformer~\cite{brehmer2023} and Clifford
neural layers~\cite{brandstetter2022,ruhe2023} likewise embed features in a Clifford
algebra and process them with the geometric/sandwich product, but target $E(n)$
equivariance for geometric data rather than associative memory and binding.
Rotary position embeddings~\cite{su2021roformer} are, as E6 makes precise, the
commuting axis-aligned special case of versor position binding; and the versor
editor of E7 belongs to the memory-based editing family (frozen model plus
external store), contrasted there with rank-one weight editing~\cite{meng2022rome}.
Our specific stance is that the attention
inner product can be the geometric product, turning attention into an
edit-friendly, order-aware binding mechanism while retaining softmax for
generalization. Within the small-model, synthetic regime studied here we have
shown, with reproducible measurements, that (i) encoding and conditioning govern
single-field capacity; (ii) sparsity restores linear capacity scaling;
(iii) the pure model generalizes by robustness and exact analogy but needs a
softmax readout to interpolate; and (iv) geometric attention recovers
interpolation, preserves analogy, and uniquely encodes order; (v) against
classical one-shot and trained baselines, the model's measured edge is storage
density and composable algebraic structure, not raw capacity or robustness;
(vi) versor position binding provably generalizes RoPE and matches it in trained
transformers once the frequency spectrum is designed; and (vii) under sequential
knowledge editing, the frozen-model versor memory surpasses a rank-one ROME
baseline on all three canonical metrics at 60 edits. Establishing
whether these advantages persist on pretrained language models, real data, and
the standard editing benchmarks is the natural next step.

\paragraph{Reproducibility.} A self-contained implementation of $\Cl(n,0)$, the
versor memory, the sparse memory, the generalization battery, the geometric
attention layer, and the JAX training code for E6--E7 (transformers, position
mechanisms, the ROME baseline, and the versor editor)---together with the JSON
logs behind every number reported here---accompanies this paper and is archived at
\href{https://doi.org/10.5281/zenodo.21752907}{doi.org/10.5281/zenodo.21752907}.

\begin{thebibliography}{99}
\small
\bibitem{vaswani2017} A.~Vaswani, N.~Shazeer, N.~Parmar, J.~Uszkoreit, L.~Jones,
A.~N.~Gomez, {\L}.~Kaiser, and I.~Polosukhin. Attention is all you need. In
\emph{Advances in Neural Information Processing Systems (NeurIPS)}, vol.~30, 2017.

\bibitem{plate1995} T.~A.~Plate. Holographic reduced representations.
\emph{IEEE Transactions on Neural Networks}, 6(3):623--641, 1995.

\bibitem{kanerva2009} P.~Kanerva. Hyperdimensional computing: An introduction to
computing in distributed representation with high-dimensional random vectors.
\emph{Cognitive Computation}, 1(2):139--159, 2009.

\bibitem{smolensky1990} P.~Smolensky. Tensor product variable binding and the
representation of symbolic structures in connectionist systems.
\emph{Artificial Intelligence}, 46(1--2):159--216, 1990.

\bibitem{kleyko2022} D.~Kleyko, D.~A.~Rachkovskij, E.~Osipov, and A.~Rahimi.
A survey on hyperdimensional computing aka vector symbolic architectures, part~I:
Models and data transformations. \emph{ACM Computing Surveys}, 55(6):130:1--130:40, 2022.

\bibitem{hestenes1984} D.~Hestenes and G.~Sobczyk. \emph{Clifford Algebra to
Geometric Calculus: A Unified Language for Mathematics and Physics}. Reidel,
Dordrecht, 1984.

\bibitem{dorst2007} L.~Dorst, D.~Fontijne, and S.~Mann. \emph{Geometric Algebra
for Computer Science: An Object-Oriented Approach to Geometry}. Morgan Kaufmann,
2007.

\bibitem{hopfield1982} J.~J.~Hopfield. Neural networks and physical systems with
emergent collective computational abilities. \emph{Proceedings of the National
Academy of Sciences}, 79(8):2554--2558, 1982.

\bibitem{ramsauer2020} H.~Ramsauer, B.~Sch\"afl, J.~Lehner, P.~Seidl, M.~Widrich,
T.~Adler, L.~Gruber, M.~Holzleitner, M.~Pavlovi\'c, G.~K.~Sandve, V.~Greiff,
D.~Kreil, M.~Kopp, G.~Klambauer, J.~Brandstetter, and S.~Hochreiter. Hopfield
networks is all you need. \emph{arXiv:2008.02217}; \emph{ICLR}, 2021.

\bibitem{brehmer2023} J.~Brehmer, P.~de~Haan, S.~Behrends, and T.~Cohen.
Geometric algebra transformer. In \emph{Advances in Neural Information Processing
Systems (NeurIPS)}, 2023. \emph{arXiv:2305.18415}.

\bibitem{brandstetter2022} J.~Brandstetter, R.~van~den~Berg, M.~Welling, and
J.~K.~Gupta. Clifford neural layers for PDE modeling. \emph{arXiv:2209.04934}, 2022.

\bibitem{ruhe2023} D.~Ruhe, J.~Brandstetter, and P.~Forr\'e. Clifford group
equivariant neural networks. In \emph{Advances in Neural Information Processing
Systems (NeurIPS)}, 2023.

\bibitem{su2021roformer} J.~Su, M.~Ahmed, Y.~Lu, S.~Pan, W.~Bo, and Y.~Liu.
RoFormer: Enhanced transformer with rotary position embedding.
\emph{Neurocomputing}, 568:127063, 2024. \emph{arXiv:2104.09864}.

\bibitem{meng2022rome} K.~Meng, D.~Bau, A.~Andonian, and Y.~Belinkov.
Locating and editing factual associations in GPT. In \emph{Advances in Neural
Information Processing Systems (NeurIPS)}, 2022. \emph{arXiv:2202.05262}.
\end{thebibliography}

\end{document}
